%%%%%%%%%%  scan idx 499  |  folio 486  %%%%%%%%%%
\SGAPrintPage{486}%
\sgaStableTarget{sga6.e08.continuation-of-the-proof-of-theorem-3-6}The proof of
Theorem 3.6 will occupy the rest of this exposé. Let us give some preliminaries.

\sgaStableTarget{sga6.e08.prop.3-7}\textbf{Proposition 3.7.} \uline{For arbitrary} $x\in \Gr K^{\bullet}(X)_{\mathbb Q}$, $y\in \Gr K^{\bullet}(Y)_{\mathbb Q}$, \uline{we have}:
\[
 f_{\mathrm{gr}}\bigl(x\cdot f^{\mathrm{gr}}(y)\bigr)
 =
 f_{\mathrm{gr}}(x)\,y .
\]
By additivity, one may suppose \(x\) and \(y\) homogeneous; say $x\in \Gr^iK^{\bullet}(X)_{\mathbb Q}$, $y\in \Gr^jK^{\bullet}(Y)_{\mathbb Q}$.
Let $x'\in \Fil^iK^{\bullet}(X)_{\mathbb Q}$, $y'\in \Fil^jK^{\bullet}(Y)_{\mathbb Q}$ be elements lifting \(x\) and \(y\). Then the projection formula gives:
\[
 f_*\bigl(x'\cdot f^*(y')\bigr)=f_*(x')\,y'.
\]
By 3.2, this equality is written in
\(\Fil^{i+j-d}K^{\bullet}(Y)_{\mathbb Q}\). It gives an equality in
\(\Gr^{i+j-d}K^{\bullet}(Y)_{\mathbb Q}\), and one checks at once that this is indeed the required equality.

\sgaStableTarget{sga6.e08.lem.3-8}\textbf{Lemma 3.8.} \uline{Let} \(f:X\to Y\), \uline{and} \(g:Y\to Z\) \uline{be two morphisms satisfying the hypotheses of} 3.6. \uline{Then} \(g\circ f\) \uline{also satisfies them}. \uline{If} \(x\) \uline{is an element of} \(K^{\bullet}(X)\), \uline{and if} 3.6 \uline{holds for} \(f\) \uline{and the element} \(x\), \uline{and for} \(g\) \uline{and the element} \(f_*(x)\), \uline{then it holds for} \(g\circ f\) \uline{and the element} \(x\).

By hypothesis, we have:
\sgaStableTarget{sga6.e08.eq.3-8-1}\begin{equation}\tag{3.8.1}
 \ch(f_*(x))=
 f_{\mathrm{gr}}\bigl(\ch(x)\cdot\Todd(\check T_f)\bigr),
\end{equation}
and:
\sgaStableTarget{sga6.e08.eq.3-8-2}\begin{equation}\tag{3.8.2}
 \ch(g_*(f_*(x)))=
 g_{\mathrm{gr}}\bigl(\ch(f_*(x))\cdot\Todd(\check T_g)\bigr).
\end{equation}

Combining (3.8.1) and (3.8.2), we obtain:
\[
\ch(g_*(f_*(x)))=
g_{\mathrm{gr}}\!\left(
 f_{\mathrm{gr}}\bigl(\ch(x)\Todd(\check T_f)\bigr)\cdot \Todd(\check T_g)
\right).
\]

%%%%%%%%%%  scan idx 500  |  folio 487  %%%%%%%%%%
\SGAPrintPage{487}%
Applying
3.7, we find:
\[
\ch(g_*(f_*(x)))=
g_{\mathrm{gr}}\!\left(
 f_{\mathrm{gr}}\bigl(\ch(x)\cdot\Todd(\check T_f)\cdot
 f^{\mathrm{gr}}(\Todd(\check T_g))\bigr)
\right).
\]
The functoriality of \(c\) and of \(\Todd\) implies:
\[
 f^{\mathrm{gr}}(\Todd(\check T_g))=\Todd(f^*(\check T_g)).
\]
By 2.7, we have:
\[
 \check T_{g\circ f}=\check T_f+f^*(\check T_g).
\]
The multiplicativity of the Todd character therefore implies:
\[
 \Todd(\check T_f)\cdot f^{\mathrm{gr}}(\Todd(\check T_g))
 =
 \Todd(\check T_{g\circ f}).
\]
Thus finally we obtain:
\[
 \ch(g_*(f_*(x)))=
 g_{\mathrm{gr}}\bigl(f_{\mathrm{gr}}\bigl(\ch(x)\cdot\Todd(\check T_{g\circ f})\bigr)\bigr),
\]
which is the announced result.

\sgaStableTarget{sga6.e08.para.3-9}\textbf{3.9.} Applying 3.8 to the factorization of \(f\) given in 3.1, one sees that, to prove the Riemann--Roch theorem, it suffices to prove it in the case of a regular immersion of constant codimension, and in the case of the structural morphism: $\mathbf{P}^r_Y\longrightarrow Y$.

Furthermore, one may suppose that the immersion \(i\) satisfies the condition of VII 4.1. Indeed, let \(s\) be the section of \(\mathbf{P}^2_{\mathbf{P}^r_Y}\) introduced in VII 4.2. Then \(s\circ i\) is a regular immersion satisfying the condition of VII 4.1 (VII 4.2), and by 3.8 it suffices to prove 3.6 for \(s\circ i\), \(\mathbf{P}^2_{\mathbf{P}^r_Y}\longrightarrow\mathbf{P}^r_Y\) and \(\mathbf{P}^r_Y\longrightarrow Y\).

We shall therefore prove 3.6 by proving the following two special cases:
\begin{enumerate}[label=\roman*)]
\item \(f\) is a regular immersion of constant codimension, satisfying the condition of VII 4.1;

%%%%%%%%%%  scan idx 501  |  folio 488  %%%%%%%%%%
\SGAPrintPage{488}%
\item \(f\) is the
structural morphism \(\mathbf{P}^r_Y\to Y\).
\end{enumerate}

\sgaStableTarget{sga6.e08.sec.4}\section*{4. Riemann--Roch theorem: case of a regular closed immersion}

\sgaStableTarget{sga6.e08.para.4-1}\textbf{4.1.} Let \(i:Y\to X\) be a regular closed immersion. The Riemann--Roch theorem for \(i\) is essentially relation VII 4.4, which makes it possible to compute the Chern classes of \(i_*(x)\) from the \(\gamma^p(N,x)\). More precisely, if one introduces the \(\lambda\)-ring \(K^{\bullet}(Y)_{\mathbb Q N}\) defined in V 5.5, then the homomorphism from \(K^{\bullet}(Y)_{\mathbb Q N}\) to \(K^{\bullet}(X)_{\mathbb Q}\) obtained from \(i_*\) is a \(\lambda\)-homomorphism by VII 2.8 and VII 4.3. The Riemann--Roch theorem then reduces to an analogous formula concerning the (additive) homomorphism \(K^{\bullet}(Y)\longrightarrow K^{\bullet}(Y)_{\mathbb Q N}\) defined by \(y\rightsquigarrow(0,y)\), that is, to a formal relation valid for every \(\lambda\)-ring. We shall not develop this proof here, in order to avoid still more formalism on \(\lambda\)-rings (cf. \(\RRR\) I 1.5), and shall return to the method of [1], which used blow-up schemes (a method already used in the proof of VII 4.3, and which is also that of the end of \(\RRR\)).

As indicated in 3.9, one may restrict to the case of a regular immersion of constant codimension \(d\) satisfying the condition of VII 4.1, the case of a general regular immersion then following from 3.8. In this section we return to the notation of VII 3.1. We thus have the cartesian square:
\[
\begin{tikzcd}[column sep=large,row sep=large]
Y' \arrow[r,"j"] \arrow[d,"g"'] & X' \arrow[d,"f"]\\
Y \arrow[r,"i"] & X ,
\end{tikzcd}
\]

%%%%%%%%%%  scan idx 502  |  folio 489  %%%%%%%%%%
\SGAPrintPage{489}%
where \(X'\) is the \(X\)-scheme
obtained by blowing up \(Y\). We suppose in addition that \(X\) is quasi-compact and has an ample sheaf.

\sgaStableTarget{sga6.e08.prop.4-2}\textbf{Proposition 4.2.} \uline{The morphism} \(f\) \uline{is projective}, \uline{a complete intersection}, \uline{and of virtual relative dimension zero}. \uline{We also have the relation}:
\[
 f_{\mathrm{gr}}\circ f^{\mathrm{gr}}=\operatorname{Id}_{\Gr K^{\bullet}(X)_{\mathbb Q}} .
\]

The first assertion follows immediately from (VII 1.8).

The second relation follows from the relation:
\[
 f_*\circ f^*=\operatorname{Id}_{K^{\bullet}(X)},
\]
proved in (VII 3.6).

\sgaStableTarget{sga6.e08.prop.4-3}\textbf{Proposition 4.3.} \uline{For every} \(y\in \Gr K^{\bullet}(Y)_{\mathbb Q}\), \uline{we have the relation}:
\[
 f^{\mathrm{gr}}\circ i_{\mathrm{gr}}(y)
 =
 j_{\mathrm{gr}}\bigl(g^{\mathrm{gr}}(y)\cdot c^{d-1}(\check F)\bigr),
\]
\uline{where} \(c^{d-1}\) \uline{denotes the} \((d-1)\)-\uline{th Chern class}.

This relation is contained in VII 4.8, since it is precisely \(v^{\mathrm{gr}}\circ u^{\mathrm{gr}}=0\).

\sgaStableTarget{sga6.e08.prop.4-4}\textbf{Proposition 4.4.} \uline{Let} \(X\) \uline{be a quasi-compact scheme}. \uline{If} \(\underline{E}\) \uline{is a locally free sheaf of finite type on} \(X\), \uline{of class} \(E\) \uline{in} \(K^{\bullet}(X)\), \uline{and of rank} \(n\), \uline{we have}:
\[
 \ch(\lambda_{-1}(E))=c^n(\check E)(\Todd(\check E))^{-1}.
\]

The two sides of the relation to be proved are multiplicative: this is clear for \(\ch\circ\lambda_{-1}\), and for \(\Todd\); it also holds for \(c^n\), which is the class of \((-1)^{\varepsilon(E)}\lambda_{-1}(E)\). Consequently, by passing to the flag bundle of \(\underline{E}\), one may suppose that \(E\) is a sum of classes of invertible sheaves, and reduce to proving the formula for an invertible sheaf. Let

%%%%%%%%%%  scan idx 503  |  folio 490  %%%%%%%%%%
\SGAPrintPage{490}%
\(\xi\) be the class of such a sheaf. We then have:
\[
 \ch(\lambda_{-1}(\xi))=\ch(1-\xi)=1-\ch(\xi).
\]
Since \(\widetilde c(\xi)=(1,1+c^1(\xi))\), the definition of \(\ch\) gives:
\[
 \ch(\lambda_{-1}(\xi))=1-\exp(c^1(\xi)).
\]

Furthermore, \(\xi\) has augmentation \(1\); the Chern class occurring in formula 4.4 for \(\xi\) is therefore \(c^1\). Taking into account the definition of \(\Todd\), and 3.5, the relation to be proved is written:
\[
 1-\exp(c^1(\xi))
 =
 -c^1(\xi)\cdot\left(\frac{-c^1(\xi)}{1-\exp(c^1(\xi))}\right)^{-1},
\]
and is therefore verified.

\sgaStableTarget{sga6.e08.rem.4-4-1}\textbf{Remark 4.4.1.} It is also possible to deduce 4.4 from a corresponding purely algebraic statement, cf. \(\RRR\) I 1.3.

\sgaStableTarget{sga6.e08.para.4-5}\textbf{4.5.} We shall now show that, to prove the Riemann--Roch theorem in the case considered in this section, one may reduce to the case where \(Y\) is a divisor of \(X\), and where the element to which the formula applies is the inverse image of an element of \(K^{\bullet}(X)\).

By 4.2, the formula to be proved:
\[
 \ch(i_*(x))=
 i_{\mathrm{gr}}\bigl(\ch(x)\cdot\Todd(\check T_i)\bigr)
\]
is equivalent to the formula:
\[
 f^{\mathrm{gr}}\bigl(\ch(i_*(x))\bigr)
 =
 f^{\mathrm{gr}}\circ i_{\mathrm{gr}}\bigl(\ch(x)\cdot\Todd(\check T_i)\bigr).
\]
Taking into account the functoriality of \(\ch\), and 4.3, it may also be written:
\[
 \ch\bigl(f^*(i_*(x))\bigr)
 =
 j_{\mathrm{gr}}\bigl(g^{\mathrm{gr}}\bigl(\ch(x)\cdot\Todd(\check T_i)\bigr)\cdot c^{d-1}(\check F)\bigr),
\]
or, by VII 3.4:
\[
 \ch\bigl(j_*(g^*(x)\cdot\lambda_{-1}(F))\bigr)
 =
 j_{\mathrm{gr}}\bigl(\ch(g^*(x))\cdot\Todd(g^*(\check T_i))\cdot c^{d-1}(\check F)\bigr).
\]
